\section{Reproducibilidad}
\label{sec:reproducibilidad}

\subsection{Estructura del repositorio}

\begin{lstlisting}
src/
  paths.py          rutas del repo, ancladas a la raiz
  data_adapter.py   carga y normaliza SST-2, AG News y Yelp
  models.py         fabrica de modelos, cabeza configurable y congelado
  train.py          bucle de fine-tuning (escrito a mano, sin Trainer)
  metrics.py        accuracy/precision/recall/F1 + latencia, memoria, tamano
  runs.py           consulta y seleccion de las corridas guardadas
  style.py          paleta y estilo comun de las figuras
  plots.py          figuras del informe (PDF + PNG)
  tables.py         tablas del informe, en Markdown
  benchmark.py      latencia y memoria en una rejilla controlada
scripts/            lanzadores de Slurm
results/
  metrics/          un JSON por corrida (nada se sobrescribe)
  figures/          figuras y tablas generadas por src/plots.py
  logs/             salida de los jobs de Slurm
docs/informe/       este documento
\end{lstlisting}

\subsection{Cómo se reproduce}

\begin{lstlisting}
bash scripts/launch_base.sh distilbert    # 3 corridas base
bash scripts/launch_base.sh bert          # 3 corridas base
CONFIGS=6 sbatch scripts/ablation.sbatch  # 18 corridas del ablation
MODELO=bert CONFIGS=mejor MEJOR="lineal|--head-hidden none" \
    sbatch scripts/ablation.sbatch        # 3 corridas de la ganadora con BERT
sbatch scripts/benchmark.sbatch           # latencia y memoria controladas
python -m src.plots                       # figuras y tablas
\end{lstlisting}

Son \num{24} corridas, unas \num{2,5} horas de GPU en una RTX A6000, en serie
porque la QOS del clúster permite una GPU a la vez.

\subsection{Trazabilidad de las cifras}

Ninguna cifra de este informe se transcribió a mano desde una consola. Cada
entrenamiento escribe un JSON con nombre único
(\texttt{\{modelo\}\_\{tarea\}[\_\{tag\}]\_\{fecha-hora\}.json}) que contiene
las métricas, la configuración completa (\texttt{vars(args)}), el historial de
\textit{loss} por iteración, el desglose por clase y la matriz de confusión.
Nada se sobrescribe nunca. Las tablas y figuras se generan desde esos archivos
con \texttt{python -m src.plots}, que produce además las versiones en Markdown
de cada tabla en \texttt{results/figures/}.

Las etiquetas (\texttt{--tag}) distinguen el propósito de cada corrida:
\texttt{base} para las del informe, \texttt{abl-<config>} para el estudio de
ablación, y \texttt{smoke} o \texttt{pilot} para pruebas y exploraciones, que
quedan excluidas automáticamente de las figuras.

\subsection{Regenerar las figuras sin GPU}

Los JSON están versionados, de modo que todas las figuras y tablas se
reconstruyen sin GPU, sin \texttt{torch} y sin volver a entrenar:

\begin{lstlisting}
pip install matplotlib numpy
python -m src.plots
\end{lstlisting}

La salida es determinista: dos ejecuciones consecutivas producen archivos
idénticos byte a byte, de modo que un archivo marcado como modificado
significa que algún dato cambió.

\subsection{Compilar este informe}

\begin{lstlisting}
cd docs/informe
make              # usa tectonic; genera main.pdf
\end{lstlisting}

Las figuras se leen de \texttt{../../results/figures/} en su versión PDF
(vectorial). El informe nunca contiene una copia de una figura, así que no
puede quedar desincronizado respecto del \textit{pipeline}.
